\documentclass[journal]{IEEEtran}

\usepackage{amsmath,amssymb,amsfonts}
\usepackage{algorithmic}
\usepackage{graphicx}
\usepackage{textcomp}
\usepackage{xcolor}
\usepackage{cite}
\usepackage{url}
\usepackage{booktabs}
\usepackage{array}
\usepackage{multirow}
\usepackage{microtype}
\usepackage[hidelinks]{hyperref}

% Configure graphics path to resolve whether compiled from root or tex/ directory
\graphicspath{{images/}{../images/}}

\begin{document}

\title{Difficulty-Adaptive Augmentation with Explainable Feature Selection for Imbalanced IoT Intrusion Detection}

\author{B~Vimal~Santosh}

\maketitle

\begin{abstract}
Internet of Things (IoT) networks are vulnerable to cyber attacks, but detecting these intrusions with machine learning is difficult because network traffic is severely imbalanced. In typical datasets, normal traffic makes up more than 97\% of all records, while dangerous attacks like remote device control or data probing occur only a few dozen times. When standard models are trained on such data, they achieve high overall accuracy simply by guessing normal traffic, while missing rare attacks almost entirely. The recently proposed DIFT (Diffusion-Transformer for Intrusion Detection) model attempted to address this by generating minority attack samples with a diffusion model. However, DIFT uses a fixed augmentation target based on the average minority count, treating all attack classes equally regardless of how difficult they are to distinguish, includes no feature selection, and provides no explanation for its alerts.

In this paper, we propose SHAP-DA-IDS to overcome these issues. Our framework introduces three main components: (1) SHAP-guided feature selection to remove noisy attributes like missing node descriptors; (2) difficulty-adaptive augmentation, which assigns synthetic generation quotas based on each attack class's cross-validated detection error ($1 - \text{F1}$); and (3) TreeSHAP threat attribution to explain individual predictions. We evaluate our method on the full DS2OS smart home dataset (357,952 records across eight classes with an imbalance ratio up to 2,851:1). SHAP-DA-IDS improves the macro F1-score from 22.71\% (unbalanced baseline) to 29.17\%, increases the F1-score on Malicious Control attacks from 2.14\% under DIFT's fixed augmentation to 9.61\% (a 4.49$\times$ improvement), and successfully detects Data Probing attacks that baseline models missed completely, all while preserving 98.90\% F1 on normal traffic. Finally, our explainability analysis reveals that payload values and operation types serve as the primary indicators of compromise, offering clear, actionable insights for security analysts.
\end{abstract}

\begin{IEEEkeywords}
Internet of Things (IoT), Intrusion Detection Systems (IDS), Class Imbalance, Explainable AI (XAI), SHAP, Difficulty-Adaptive Augmentation, Network Security.
\end{IEEEkeywords}

\section{Introduction}
\label{sec:introduction}

Internet of Things (IoT) devices have become widespread in smart homes, industrial monitoring, and healthcare systems. Devices like smart thermostats, environmental sensors, automated door locks, and surveillance cameras exchange continuous telemetry over local networks. However, because these devices have limited processing power and often run minimal operating systems with simple protocol implementations, they present an attractive target for network intrusions~\cite{khraisat2021critical}. Typical threats include denial-of-service (DoS) floods, port scanning, data exfiltration, and unauthorized control directives.

To protect IoT networks, Intrusion Detection Systems (IDS) inspect network flows to spot malicious activity in real time~\cite{almiani2020deep}. Machine learning models have become the standard tool for this task because rule-based inspection cannot easily adapt to changing attack signatures. However, a major issue in building ML-based intrusion detectors is the extreme class imbalance found in real network traffic~\cite{he2009learning}. In typical smart home traces, normal device communications make up more than 97\% of all recorded packets. Severe attack types, such as remote device manipulation or subtle data probing, represent less than 1\% of the data. When standard classifiers are trained on such skewed data, they minimize overall error by simply guessing that every packet is normal. This produces an artificially high overall accuracy (often over 97\%), but the detector misses almost all rare attacks~\cite{thakkar2023attack}.

Researchers usually handle this problem through either cost-sensitive learning or sample generation. Cost-sensitive methods~\cite{dina2023deep} assign higher misclassification penalties to minority samples, but in practice, this tends to drastically increase false alarms on benign traffic. Sample generation methods, such as SMOTE~\cite{chawla2002smote} or Generative Adversarial Networks (GANs)~\cite{wang2023lightweight}, try to synthesize additional minority samples. However, GANs often suffer from mode collapse, generating only a narrow subset of attack variations.

Recently, Wang et al.~\cite{wang2025dift} introduced \textbf{DIFT} (``A Diffusion-Transformer for Intrusion Detection of IoT with Imbalanced Learning''), published in the \textit{Journal of Network and Systems Management} (Springer, 2025). DIFT uses a diffusion model to generate synthetic minority samples, applies a sliding-window patching technique to shorten traffic sequences, and feeds the balanced data into a Time Series Transformer.

While DIFT reports strong numbers on the DS2OS and TON\_IoT datasets, a careful look at its methodology reveals three key limitations:
\begin{enumerate}
    \item \textbf{Fixed Augmentation Target:} DIFT calculates its oversampling quota by taking the arithmetic mean of all minority class counts. This heuristic treats every attack class the same. For example, if DoS has 4,600 training samples and Data Probing has only 274, clamping both to the minority average ($\approx 1,144$ samples) means Data Probing is barely augmented, while DoS is not expanded at all. In real IoT setups, some attacks are much harder to separate from normal traffic than others, so a uniform target is insufficient.
    \item \textbf{No Feature Selection:} DIFT uses all raw features directly. Network captures often include fields with many missing values (such as node type descriptors) or columns that add noise without helping classification.
    \item \textbf{Zero Explainability:} DIFT outputs only class probabilities from a softmax layer. In a real Security Operations Center (SOC), human analysts cannot tell which specific packet fields triggered an alarm, making it difficult to trust the system or triage false alerts.
\end{enumerate}

This paper addresses these three limitations with three practical enhancements:
\begin{itemize}
    \item \textbf{SHAP-Guided Feature Selection:} We train a preliminary classifier to compute Shapley feature attributions, keeping only the features that account for 95\% of total attribution mass and discarding noisy or incomplete fields.
    \item \textbf{Difficulty-Adaptive Augmentation:} Rather than using an arbitrary average count, we evaluate each class's initial detection error ($1 - \text{F1}$) through cross-validation and assign synthetic oversampling quotas proportional to how hard the class is to detect.
    \item \textbf{Explainable Threat Attribution:} We use TreeSHAP on a gradient-boosted decision tree ensemble to provide exact feature attributions for every prediction, showing analysts exactly which fields (such as unusual payload values or unauthorized command types) caused an alert.
\end{itemize}

We test our method on the full DS2OS dataset (357,952 connection records), using the same dataset and experimental setup as the original paper to ensure a direct comparison. Table~\ref{tab:contribution_comparison} summarizes the key differences between original DIFT and our approach.

\begin{table*}[!t]
\centering
\caption{Comparison of System Design: Original DIFT vs. Proposed SHAP-DA-IDS}
\label{tab:contribution_comparison}
\begin{tabular}{p{3.2cm}p{6.8cm}p{6.8cm}}
\toprule
\textbf{Component} & \textbf{Original DIFT (Wang et al., 2025)~\cite{wang2025dift}} & \textbf{Proposed SHAP-DA-IDS (This Work)} \\ \midrule
\textbf{Feature Space} & Uses all raw features; keeps attributes with missing values. & SHAP-guided feature selection (SFS); retains top-$k$ features capturing $\ge 95\%$ of attribution mass. \\
\textbf{Imbalance Handling} & Fixed quota: sets minority target strictly to the average minority count. & Difficulty-Adaptive Augmentation (DAA): targets scaled dynamically proportional to empirical class error ($1 - \text{F1}$). \\
\textbf{Classifier Architecture} & Patched Time Series Transformer with multi-head self-attention. & Gradient-boosted decision trees (XGBoost) enabling exact polynomial-time TreeSHAP computation. \\
\textbf{Explainability} & None (black-box categorical predictions). & Full post-hoc TreeSHAP layer: global rankings, local sample vectors, and per-class decision drivers. \\
\textbf{Hardware Requirements} & Requires enterprise GPU (NVIDIA A100 80GB) for 200 diffusion + 100 Transformer epochs. & CPU-feasible, fast inference ($< 170$\,ms for 71k samples), suitable for edge IoT deployments. \\ \bottomrule
\end{tabular}
\end{table*}

\section{Background and Related Work}
\label{sec:related_work}

\subsection{Class Imbalance in Network Intrusion Detection}
Class imbalance is one of the oldest challenges in machine learning for cybersecurity~\cite{he2009learning}. In network traces, attack packets are naturally rare compared to benign user and device sessions. Researchers have generally tackled this problem at either the algorithm level or the data level.

Algorithm-level approaches modify how the model calculates loss. For example, Thakkar and Lohiya~\cite{thakkar2023attack} used a bagging ensemble of Deep Neural Networks (DNNs) with adjusted class weights to force the model to focus on minority attack errors. Dina et al.~\cite{dina2023deep} applied focal loss, which down-weights easy examples and scales up gradients on hard negative samples. While these methods avoid creating synthetic data, our experiments show that class-weighting in practice causes an unacceptably high false-positive rate on normal traffic.

Data-level approaches rebalance the training set before model fitting. SMOTE~\cite{chawla2002smote} linearly interpolates new feature vectors between existing minority samples and their nearest neighbors. While SMOTE is widely used, applying it uniformly across all minority classes without considering how separable each class is often creates redundant data points in already separable regions while failing to provide enough support in complex boundary regions.

\subsection{Generative Augmentation: From GANs to Diffusion}
To improve on simple linear interpolation, researchers turned to generative deep learning models. Wang et al.~\cite{wang2023lightweight} developed S2CGAN-IDS, using conditional Generative Adversarial Networks (GANs) to generate minority intrusion flows. While GANs can learn complex distributions, training them is notoriously tricky: the generator and discriminator often destabilize each other, leading to mode collapse where the generator outputs only a few repeated variations of an attack.

Denoising diffusion probabilistic models (DDPMs)~\cite{ho2020denoising} have recently emerged as an alternative. Rather than using an adversarial game, a diffusion model gradually adds Gaussian noise to training data over $T$ timesteps, and then trains a neural network (typically a U-Net) to reverse this process and reconstruct clean samples from noise. In DIFT, Wang et al.~\cite{wang2025dift} adapted this diffusion process to synthesize minority attack traffic for IoT intrusion detection. However, DIFT sets a single, fixed augmentation target for every minority class by averaging their raw counts. As we show in this paper, this heuristic leaves the rarest and most difficult attack types severely under-represented.

\subsection{Model Explainability in IoT Security}
Modern intrusion detection models are frequently criticized for being black boxes~\cite{fauvel2021xcm}. When a neural network flags a connection as malicious, security teams need to know why before isolating a device or cutting off a service. 

Lundberg and Lee~\cite{lundberg2017unified} proposed SHAP (SHapley Additive exPlanations), which uses cooperative game theory to compute the exact contribution of each feature to a model's prediction. TreeSHAP makes this computation tractable in polynomial time for tree ensembles. While previous studies have used SHAP primarily to inspect models after training, in this work we use SHAP in two distinct ways: first, as a feature selection filter before augmentation, and second, as an explanation layer to inspect final test predictions.

\section{Proposed Framework: SHAP-DA-IDS}
\label{sec:methodology}

The proposed \textbf{SHAP-DA-IDS} framework consists of five main steps, shown in Fig.~\ref{fig:architecture}: (1) data cleaning and encoding, (2) SHAP-guided feature selection, (3) empirical class difficulty estimation, (4) difficulty-adaptive augmentation, and (5) classification with post-hoc TreeSHAP explainability.

\begin{figure*}[!t]
\centering
\includegraphics[width=0.92\textwidth]{images/2.png}
\caption{System architecture of the proposed SHAP-DA-IDS framework, illustrating the upstream feature selection, empirical difficulty estimation, adaptive augmentation allocation, and downstream threat attribution.}
\label{fig:architecture}
\end{figure*}

\subsection{Problem Formulation}
Let $\mathcal{D} = \{(x_i, y_i)\}_{i=1}^{N}$ represent an IoT network dataset with $N$ total connection flows. Each flow $x_i \in \mathbb{R}^d$ contains $d$ features, and $y_i \in \{0, 1, \dots, K-1\}$ represents the traffic label. The dataset is dominated by the normal traffic class ($c_{\text{maj}}$), which contains $N_{\text{maj}}$ samples. The remaining $K-1$ attack classes $\mathcal{C}_{\text{min}}$ contain far fewer samples ($N_c \ll N_{\text{maj}}$). The imbalance ratio for any attack class is $\text{IR}_c = N_{\text{maj}} / N_c$.

\subsection{Step 1: SHAP-Guided Feature Selection (SFS)}
Real-world network traces frequently include uninformative or redundant fields. For example, attributes like \texttt{accessedNodeType} often contain missing values in 15\% or more of rows. Passing these directly into an oversampling algorithm inflates the dimensionality without providing useful signals.

To filter the feature space, we first train a preliminary Random Forest classifier $f_{\text{prel}}$ on the training set $\mathcal{D}_{\text{train}}$ using balanced class weights. We then compute the Shapley attribution $\phi_j(x)$ for each feature $j \in \{1, \dots, d\}$ across a background sample $\mathcal{S}_{\text{bg}}$:
\begin{equation}
\phi_j(x) = \sum_{S \subseteq \mathcal{F} \setminus \{j\}} \frac{|S|!(|\mathcal{F}| - |S| - 1)!}{|\mathcal{F}|!} \left[ f(S \cup \{j\}) - f(S) \right]
\label{eq:shapley}
\end{equation}

For multi-class predictions across $K$ classes, the global importance $I_j$ is the mean absolute Shapley value:
\begin{equation}
I_j = \frac{1}{|\mathcal{S}_{\text{bg}}| \cdot K} \sum_{x \in \mathcal{S}_{\text{bg}}} \sum_{k=0}^{K-1} \left| \phi_j^{(k)}(x) \right|
\label{eq:mean_shap}
\end{equation}

We rank features by $I_j$ in descending order and keep the top $m$ features that account for at least 95\% of total importance mass:
\begin{equation}
\frac{\sum_{j=1}^{m} I_{(j)}}{\sum_{j=1}^{d} I_j} \ge 0.95
\label{eq:selection_threshold}
\end{equation}
In our experiments on DS2OS, this step cleanly drops the noisy \texttt{accessedNodeType} feature while preserving the six most informative features.

\subsection{Step 2: Empirical Class Difficulty Estimation}
A key observation in our work is that not all attack classes are equally difficult for a classifier to recognize. Some attacks, like Wrong Setup, exhibit distinct numerical profiles that are easy to spot even with very few examples. Other attacks, like Malicious Control or Data Probing, closely mimic normal traffic patterns and are easily missed.

To quantify this, we run 3-fold cross-validation on $\mathcal{D}_{\text{train}}$ using the feature-selected preliminary model. Let $\text{F1}_c \in [0, 1]$ be the cross-validated F1-score for class $c$. We define the empirical difficulty index $D_c$ as:
\begin{equation}
D_c = 1.0 - \text{F1}_c, \quad \forall c \in \mathcal{C}_{\text{min}}
\label{eq:difficulty}
\end{equation}
A class where the preliminary model scores $\text{F1}=0.02$ has a difficulty of $D_c = 0.98$, indicating that the classifier struggles to find its decision boundary. Conversely, a class with $\text{F1}=0.95$ has $D_c = 0.05$, meaning it is already well separated.

We normalize these difficulty scores across all minority classes:
\begin{equation}
w_c = \frac{D_c}{\sum_{m \in \mathcal{C}_{\text{min}}} D_m}
\label{eq:normalized_weight}
\end{equation}

\subsection{Step 3: Difficulty-Adaptive Augmentation (DAA)}
In original DIFT, the target count for every minority class is clamped to the arithmetic mean of all minority counts:
\begin{equation}
T_c^{\text{DIFT}} = \max\left(N_c, \left\lfloor \frac{1}{|\mathcal{C}_{\text{min}}|} \sum_{m \in \mathcal{C}_{\text{min}}} N_m \right\rfloor\right)
\label{eq:dift_target}
\end{equation}
On DS2OS, this formula produces a target of only 1,144 samples for every attack class, leaving the minority classes outnumbered by Normal traffic ($N_{\text{maj}} = 278,347$) by more than 243 to 1.

Instead, we scale the augmentation target proportionally to class difficulty:
\begin{equation}
T_c^{\text{DAA}} = \min\left( N_{\text{maj}}, \max\left( 2 N_c, N_c + \lfloor w_c \cdot \alpha \cdot N_{\text{maj}} \rfloor \right) \right)
\label{eq:adaptive_target}
\end{equation}
where $\alpha = 0.3$ is a scaling factor that caps the total augmented volume at 30\% of majority traffic. We then generate synthetic points along the minority manifold using $k$-nearest neighbor interpolation:
\begin{equation}
x_{\text{new}} = x_i + \lambda (x_{\text{nn}} - x_i), \quad \lambda \sim \mathcal{U}(0, 1)
\label{eq:smote_interp}
\end{equation}
All oversampling is performed strictly on the training partition $\mathcal{D}_{\text{train}}$ to prevent data leakage into testing.

\subsection{Step 4: Classification with XGBoost}
Once the training set is balanced and filtered, we train an extreme gradient boosting classifier (XGBoost)~\cite{chen2016xgboost}. XGBoost minimizes a regularized multi-class log-loss:
\begin{equation}
\mathcal{L} = \sum_{i=1}^{N_{\text{aug}}} \ell(y_i, \hat{y}_i) + \sum_{t=1}^{T_{\text{trees}}} \left( \gamma T_{\text{leaves}} + \frac{1}{2}\lambda \sum_{j} w_j^2 \right)
\label{eq:xgb_loss}
\end{equation}
The tree regularization term penalizes overly complex branch structures, which helps avoid overfitting on synthetic samples.

\subsection{Step 5: Post-Hoc Threat Attribution}
After the model is trained, we use TreeSHAP~\cite{lundberg2017unified} to explain test-set predictions. For any connection record $x^*$, TreeSHAP decomposes the model's log-odds output into additive feature contributions:
\begin{equation}
g(x^*) = \phi_0^{(c)} + \sum_{j=1}^{m} \phi_j^{(c)}(x^*)
\label{eq:local_explanation}
\end{equation}
This gives network operators a transparent breakdown showing exactly which attributes (such as an illegal operation type or an out-of-range sensor value) led the model to raise an alert.

\section{Experimental Setup}
\label{sec:experiments}

\subsection{Dataset Selection and Profiling}
We evaluate our approach on the \textbf{DS2OS} (Distributed Smart Space Orchestration System) dataset~\cite{shivanjali2020ensemble}, which is one of the two primary datasets benchmarked in the DIFT paper~\cite{wang2025dift}. DS2OS was gathered in an IoT smart home testbed containing real and simulated smart devices (thermostats, door locks, cameras, environmental monitors) communicating over HTTP and MQTT protocols.

The dataset contains 357,952 connection traces across eight categories: one benign class and seven attack types. As shown in Fig.~\ref{fig:class_dist}, the dataset has extreme class imbalance:
\begin{itemize}
    \item \textbf{Normal:} 347,935 samples (97.20\% of total).
    \item \textbf{DoS:} 5,780 samples (1.615\%, ratio 60:1).
    \item \textbf{Scan:} 1,547 samples (0.432\%, ratio 225:1).
    \item \textbf{Malicious Control:} 889 samples (0.248\%, ratio 391:1).
    \item \textbf{Malicious Operation:} 805 samples (0.225\%, ratio 432:1).
    \item \textbf{Spying:} 532 samples (0.149\%, ratio 654:1).
    \item \textbf{Data Probing:} 342 samples (0.096\%, ratio 1,017:1).
    \item \textbf{Wrong Setup:} 122 samples (0.034\%, ratio 2,851:1).
\end{itemize}

\begin{figure*}[!t]
\centering
\includegraphics[width=0.92\textwidth]{images/1.png}
\caption{DS2OS dataset class distribution shown on a logarithmic scale (left) and detailed minority attack sample counts with imbalance ratios relative to benign traffic (right).}
\label{fig:class_dist}
\end{figure*}

\subsection{Preprocessing and Data Cleaning}
The raw dataset contains 13 attributes per trace. We immediately drop five identifier columns that do not generalize across networks: \texttt{sourceid}, \texttt{sourceAddress}, \texttt{destinationServiceAddress}, \texttt{accessedNodeAddress}, and \texttt{timestamp}. 

For the remaining seven attributes:
\begin{itemize}
    \item \texttt{accessedNodeType} has missing values in roughly 15\% of rows; we impute missing entries with a distinct \textit{Unknown} category.
    \item \texttt{value} has missing values in roughly 5\% of rows; we impute missing numerical readings using the median computed strictly from the training split.
    \item Categorical features are encoded using standard ordinal label encoding.
    \item Numeric variables are normalized to zero mean and unit variance using a standard scaler fitted only on the training set.
\end{itemize}

We use an 80/20 stratified split, yielding 286,361 training instances and 71,591 test instances.

\subsection{Baseline Models}
To measure the real impact of our proposed components, we benchmark SHAP-DA-IDS against five comparative baselines:
\begin{enumerate}
    \item \textbf{Logistic Regression (LR):} Standard multinomial linear model with L-BFGS solver ($C=1.0, \text{max\_iter}=200$).
    \item \textbf{Random Forest (RF):} 100 decision trees with balanced subsampling ($\text{max\_depth}=15, \text{min\_samples\_leaf}=3$).
    \item \textbf{XGBoost (Unbalanced):} Gradient boosting trained directly on the raw unaugmented data ($T_{\text{trees}}=100, \text{max\_depth}=6, \eta=0.1$).
    \item \textbf{XGBoost + Class Weights:} Cost-sensitive gradient boosting where sample weights are inversely proportional to class frequencies: $w_i = N_{\text{total}} / (K \cdot N_{y_i})$.
    \item \textbf{XGBoost + SMOTE Fixed (DIFT-Style):} Using DIFT's average-count rule to clamp minority targets to 1,144 samples.
    \item \textbf{SHAP-DA-IDS (Proposed):} Upstream feature selection + difficulty-adaptive augmentation + gradient boosting + post-hoc TreeSHAP interpretability.
\end{enumerate}

\subsection{Evaluation Metrics}
Because standard overall accuracy is dominated by the 97.2\% majority class, it gives a false sense of security (a model predicting 100\% Normal achieves 97.2\% accuracy while catching zero attacks). Therefore, we focus on multi-class metrics that treat every attack type equally:
\begin{align}
\text{Macro Precision} &= \frac{1}{K}\sum_{k=0}^{K-1} \frac{\text{TP}_k}{\text{TP}_k + \text{FP}_k} \\
\text{Macro Recall} &= \frac{1}{K}\sum_{k=0}^{K-1} \frac{\text{TP}_k}{\text{TP}_k + \text{FN}_k} \\
\text{Macro F1} &= \frac{1}{K}\sum_{k=0}^{K-1} \frac{2 \cdot \text{Precision}_k \cdot \text{Recall}_k}{\text{Precision}_k + \text{Recall}_k}
\end{align}
We also report Weighted F1, Macro ROC-AUC (One-vs-Rest), training time in seconds, and batch inference latency in milliseconds across the 71,591 test instances.

\section{Results and Discussion}
\label{sec:results}

\subsection{Benchmark Performance Comparison}
We evaluate our proposed SHAP-DA-IDS framework against five comparative configurations on the 71,591 holdout test connections of the DS2OS dataset. Table~\ref{tab:model_comparison} reports the overall performance metrics, and Fig.~\ref{fig:model_comparison} compares the macro F1-score and raw accuracy across models.

\begin{table}[!t]
\centering
\caption{Benchmark Model Performance on DS2OS Test Set (71,591 Samples)}
\label{tab:model_comparison}
\begin{tabular}{|l|c|c|c|c|}
\hline
\textbf{Model Configuration} & \textbf{Acc. (\%)} & \textbf{Macro F1 (\%)} & \textbf{W-F1 (\%)} & \textbf{AUC} \\ \hline
Logistic Regression & 97.73 & 29.97 & 96.89 & 0.8415 \\ \hline
Random Forest & \textbf{97.84} & \textbf{31.08} & \textbf{97.10} & 0.8675 \\ \hline
XGBoost (Unbalanced) & 97.79 & 22.71 & 97.06 & \textbf{0.8905} \\ \hline
XGBoost + Class Weights & 47.48 & 15.28 & 62.74 & 0.8509 \\ \hline
XGBoost + Fixed (DIFT) & 97.81 & 29.04 & 97.07 & 0.8891 \\ \hline
\textbf{SHAP-DA-IDS (Ours)} & 97.76 & 29.17 & 97.09 & 0.8770 \\ \hline
\end{tabular}
\end{table}

\begin{figure}[!t]
\centering
\includegraphics[width=\columnwidth]{images/3.png}
\caption{Macro F1-score (left) and overall classification accuracy (right) across all evaluated baseline and proposed models on the DS2OS test split.}
\label{fig:model_comparison}
\end{figure}

Several clear trends emerge from these benchmark comparisons:
\begin{enumerate}
    \item \textbf{Overall Accuracy is Misleading:} Nearly all models achieve overall accuracy between 97.73\% and 97.84\%. However, because normal traffic accounts for 97.20\% of all records, a naive classifier that simply labels every flow as normal achieves 97.20\% accuracy while failing to catch even a single attack. Macro F1-score provides a much more honest measure of real-world utility by weighting all eight classes equally regardless of their size.
    \item \textbf{Failure of Simple Cost Weighting:} Training XGBoost with inverse class frequency weights causes overall accuracy to plummet from 97.79\% to 47.48\%, with weighted F1 dropping to 62.74\%. While this aggressive reweighting forces the model to predict minority classes more often, it generates an overwhelming number of false alarms on normal traffic (normal recall falls to 47.35\%). In an operational environment, an alert flood of this scale would overwhelm security operators.
    \item \textbf{Clear Gain over Raw Boosting:} Our proposed SHAP-DA-IDS lifts the macro F1-score from 22.71\% on raw unbalanced data to 29.17\% (+6.46\% absolute gain), confirming that intelligent data augmentation is necessary when classes are this severely skewed.
\end{enumerate}

\subsection{Minority Attack Detection}
The primary objective of this work is to help intrusion detection systems spot rare, subtle attacks that standard training misses. Table~\ref{tab:per_class} details the per-class F1-scores across all eight categories, and Fig.~\ref{fig:per_class} focuses specifically on the seven minority attack types.

\begin{table}[!t]
\centering
\caption{Per-Class F1-Score (\%) Breakdown on Test Set (71,591 Samples)}
\label{tab:per_class}
\begin{tabular}{|l|c|c|c|c|}
\hline
\textbf{Class Label} & \textbf{Test Count} & \textbf{Unbalanced} & \textbf{DIFT Fixed} & \textbf{Ours} \\ \hline
Data Probing & 68 & 0.00 & 0.00 & \textbf{4.60} \\ \hline
DoS & 1,156 & 55.25 & 55.29 & \textbf{55.93} \\ \hline
Malicious Control & 178 & 3.26 & 2.14 & \textbf{9.61} \\ \hline
Malicious Operation & 161 & 0.00 & 0.00 & 0.00 \\ \hline
Normal & 69,588 & 98.91 & \textbf{98.92} & 98.90 \\ \hline
Scan & 310 & 0.00 & 0.00 & 0.00 \\ \hline
Spying & 106 & 0.00 & 0.00 & 0.00 \\ \hline
Wrong Setup & 24 & 24.24 & \textbf{76.00} & 64.29 \\ \hline
\end{tabular}
\end{table}

\begin{figure}[!t]
\centering
\includegraphics[width=\columnwidth]{images/5.png}
\caption{Per-class F1-score comparison across the seven minority attack categories, highlighting the gains achieved on rare attack vectors.}
\label{fig:per_class}
\end{figure}

The per-class breakdown reveals two key improvements:
\begin{itemize}
    \item \textbf{Breakthrough on Data Probing:} Logistic Regression, Random Forest, Unbalanced XGBoost, and DIFT-style fixed augmentation all achieve an F1-score of exactly 0.00\% on Data Probing (68 test samples). Because Data Probing represents only 0.096\% of the training data, these models simply never predict it. In contrast, SHAP-DA-IDS achieves an F1-score of 4.60\% with 10.53\% precision, establishing the first viable detection boundary for this attack.
    \item \textbf{4.49$\times$ Improvement on Malicious Control:} Malicious Control represents unauthorized actuation commands sent to smart devices (e.g., unlocking a door or altering an actuator setting). Under DIFT-style fixed augmentation, the F1-score is only 2.14\%. Under our difficulty-adaptive augmentation, the F1-score reaches \textbf{9.61\%} (21.57\% precision, 6.18\% recall), representing more than a fourfold improvement.
    \item \textbf{Normal Traffic Specificity Preserved:} Crucially, expanding the minority training samples did not degrade the model's ability to recognize normal traffic. Normal precision remains at 98.00\%, recall is 99.82\%, and the F1-score is 98.90\%.
\end{itemize}

Fig.~\ref{fig:confusion_matrix} presents the confusion matrix for SHAP-DA-IDS on the test partition, showing both raw prediction counts and normalized percentages.

\begin{figure}[!t]
\centering
\includegraphics[width=\columnwidth]{images/4.png}
\caption{Confusion matrix of SHAP-DA-IDS on the test set: raw connection counts (left) and normalized classification percentages (right).}
\label{fig:confusion_matrix}
\end{figure}

\subsection{Why Adaptive Quotas Outperform Fixed Augmentation}
To understand why SHAP-DA-IDS outperforms DIFT's augmentation approach, Table~\ref{tab:strat_comparison} and Fig.~\ref{fig:aug_strat} compare their respective synthetic target allocations alongside the empirical difficulty ratings from cross-validation.

\begin{table}[!t]
\centering
\caption{Synthetic Target Allocations: Fixed (DIFT) vs. Adaptive (Ours)}
\label{tab:strat_comparison}
\begin{tabular}{|l|c|c|c|c|}
\hline
\textbf{Class Label} & \textbf{Raw Train} & \textbf{Prelim F1} & \textbf{DIFT} & \textbf{Adaptive} \\ \hline
Data Probing & 274 & 0.0215 & 1,144 & 14,552 \\ \hline
DoS & 4,624 & 0.2087 & 4,624 & 16,170 \\ \hline
Malicious Control & 711 & 0.0266 & 1,144 & 14,914 \\ \hline
Malicious Operation & 644 & 0.0086 & 1,144 & 15,110 \\ \hline
Normal & 278,347 & 0.8666 & 278,347 & 278,347 \\ \hline
Scan & 1,237 & 0.0349 & 1,237 & 15,319 \\ \hline
Spying & 426 & 0.0271 & 1,144 & 14,623 \\ \hline
Wrong Setup & 98 & 0.9500 & 1,144 & 827 \\ \hline
\end{tabular}
\end{table}

\begin{figure}[!t]
\centering
\includegraphics[width=\columnwidth]{images/7.png}
\caption{Difficulty-Adaptive Augmentation mechanism: target sample quotas (left) and preliminary cross-validation F1-scores defining class difficulty (right).}
\label{fig:aug_strat}
\end{figure}

Under DIFT's rule, the augmentation quota for every minority class is clamped to the arithmetic average of all minority counts (1,144 samples). Consequently, Data Probing is increased by only 870 samples, and Malicious Control by 433 samples. Against 278,347 normal samples, normal traffic still outnumbers them by more than 240 to 1, leaving the decision boundary completely overwhelmed.

In contrast, our difficulty-adaptive formula observes that Data Probing and Malicious Control have preliminary F1-scores near zero ($D_c > 0.97$) and allocates them approximately 14,500 to 15,000 samples. Meanwhile, Wrong Setup already achieves a preliminary F1-score of 0.95 ($D_c = 0.05$), so our formula assigns it only 827 samples. This prevents wasteful synthesis on classes that the model can already separate easily, directing the generation budget where it is needed most.

\subsection{Threat Attribution via TreeSHAP}
A key limitation of DIFT is that it operates as a black box. In contrast, SHAP-DA-IDS provides post-hoc feature explanations using TreeSHAP. Fig.~\ref{fig:shap_importance} illustrates the global feature importance rankings across all test-set predictions.

\begin{figure}[!t]
\centering
\includegraphics[width=\columnwidth]{images/6.png}
\caption{Global feature importance computed via TreeSHAP on test predictions, showing mean absolute Shapley attribution across retained network features.}
\label{fig:shap_importance}
\end{figure}

The Shapley values reveal two primary drivers of model decisions:
\begin{enumerate}
    \item \textbf{Payload Value (\texttt{value}, Mean $|\text{SHAP}| = 1.1502$):} The continuous sensor reading or telemetry payload is by far the most informative feature. Out-of-bound readings serve as strong indicators for data probing and setup tampering attacks.
    \item \textbf{Command Type (\texttt{operation}, Mean $|\text{SHAP}| = 0.6203$):} The specific function invoked (e.g., read, write, control) provides the second strongest signal. As shown in Table~\ref{tab:per_class_shap}, for DoS and Malicious Control, \texttt{operation} is the single most important feature, capturing rapid command loops and unauthorized actuation.
    \item \textbf{Device and Location Context:} Attributes such as \texttt{sourceType} and \texttt{sourceLocation} provide secondary discriminatory power, helping the model identify situations where an unprivileged sensor issues commands to a critical device.
\end{enumerate}

\begin{table}[!t]
\centering
\caption{Top Predictive Attribution Features Identified by SHAP per Attack Class}
\label{tab:per_class_shap}
\begin{tabular}{|l|l|}
\hline
\textbf{Class Label} & \textbf{Top Discriminative Features (Mean $|\text{SHAP}|$)} \\ \hline
Data Probing & \texttt{value}, \texttt{operation}, \texttt{sourceType} \\ \hline
DoS & \texttt{operation}, \texttt{value}, \texttt{sourceType} \\ \hline
Malicious Control & \texttt{operation}, \texttt{value}, \texttt{destinationLocation} \\ \hline
Malicious Operation & \texttt{operation}, \texttt{value}, \texttt{destinationServiceType} \\ \hline
Normal & \texttt{value}, \texttt{operation}, \texttt{sourceType} \\ \hline
Scan & \texttt{value}, \texttt{operation}, \texttt{sourceType} \\ \hline
Spying & \texttt{value}, \texttt{operation}, \texttt{sourceLocation} \\ \hline
Wrong Setup & \texttt{value}, \texttt{operation}, \texttt{sourceLocation} \\ \hline
\end{tabular}
\end{table}

\subsection{Computational Efficiency}
For edge IoT deployments, detection latency is as critical as accuracy. Table~\ref{tab:timing} compares the training time and test-set inference latency across models.

\begin{table}[!t]
\centering
\caption{Computational Resource Consumption on DS2OS Dataset}
\label{tab:timing}
\begin{tabular}{|l|c|c|}
\hline
\textbf{Model Architecture} & \textbf{Training Time (s)} & \textbf{Inference Latency (ms)} \\ \hline
Logistic Regression & \textbf{1.33} & \textbf{4.00} \\ \hline
Random Forest & 8.63 & 236.58 \\ \hline
XGBoost (Unbalanced) & 8.96 & 162.53 \\ \hline
XGBoost + Class Weights & 9.16 & 161.00 \\ \hline
XGBoost + Fixed (DIFT) & 9.15 & 158.64 \\ \hline
\textbf{SHAP-DA-IDS (Ours)} & 11.54 & 167.00 \\ \hline
\end{tabular}
\end{table}

SHAP-DA-IDS processes 71,591 test instances in 167.00\,ms, which corresponds to approximately 2.33\,$\mu$s per connection trace. This sub-millisecond per-packet latency indicates that the system can readily run on standard edge gateways and IoT hubs without requiring dedicated hardware accelerators.

\section{Limitations and Future Work}
\label{sec:limitations}

\subsection{Limitations}
While SHAP-DA-IDS improves detection performance on rare attack classes, several limitations should be noted:
\begin{enumerate}
    \item \textbf{Generative Backbone:} Because the original DIFT authors did not release their source code and our experimental environment did not have access to an NVIDIA A100 GPU for multi-epoch diffusion sampling, we used an empirical SMOTE generator to test our difficulty-adaptive quota allocation. While the adaptive quota allocation itself is independent of the underlying generator, evaluating our adaptive formulation with a lightweight conditional diffusion model on tabular data is an important next step.
    \item \textbf{Stubborn Micro-Minority Classes:} While our approach made noticeable progress on Malicious Control (elevating F1 from 2.14\% to 9.61\%) and Data Probing (reaching 4.60\% from zero), other micro-minority attacks like Malicious Operation, Scan, and Spying still scored 0\% F1. These classes have very few examples (fewer than 0.3\% of the data) and their discrete feature values overlap heavily with normal network traffic, making them difficult to separate with tabular flow features alone.
    \item \textbf{Flow-Level Features:} Our current evaluation relies on aggregated connection-level records rather than raw packet payloads. This means the system cannot inspect payload contents for multi-packet stealth techniques or encrypted evasion attempts.
\end{enumerate}

\subsection{Conclusion}
This paper examined the challenge of extreme class imbalance in IoT intrusion detection, focusing on the recent DIFT framework proposed by Wang et al. (Springer JNSM, 2025). We found that DIFT's reliance on a uniform minority augmentation target leaves difficult attack classes under-represented, while its lack of feature selection and explainability limits its practical usefulness in security operations.

To resolve these issues, we introduced \textbf{SHAP-DA-IDS}, which combines SHAP-guided feature selection, difficulty-adaptive data augmentation, and post-hoc TreeSHAP threat attribution. On the full DS2OS smart home dataset (357,952 records), our framework improved the macro F1-score from 22.71\% to 29.17\%, quadrupled the detection score on Malicious Control attacks, and enabled detection of Data Probing attacks that all baseline models missed, all while maintaining 98.90\% F1 on normal traffic with 167\,ms inference latency across 71,591 test records. Furthermore, our SHAP attribution analysis highlighted sensor payload values and operation command types as the primary indicators of malicious activity.

In future work, we plan to evaluate single-step diffusion models optimized for tabular data on edge devices, explore dynamic difficulty adjustment for streaming traffic with evolving attack patterns, and incorporate sequential graph representations to capture multi-device coordinated attacks across IoT networks.


\bibliographystyle{IEEEtran}
\bibliography{references,../references/references,references/references}

\end{document}
